\section{Background}
\label{sec:background}

\subsection{Terminal Maneuvering Area Operations}
The Terminal Maneuvering Area (TMA) represents one of the most critical and complex sectors of modern airspace. Operations within this region are characterized by high traffic density, reduced separation minimums, low altitudes, and limited runway capacity. These constraints necessitate frequent air traffic controller interventions to maintain safety and sequence arriving aircraft. However, such tactical interventions often disrupt optimal flight profiles. When an aircraft is unable to perform a Continuous Descent Operation (CDO) and is instead forced to level off, enter holding patterns, or follow extended vectoring, it operates at lower, less aerodynamically efficient altitudes. Consequently, these deviations from optimal trajectories generate substantial increases in fuel consumption and greenhouse gas emissions.

\subsection{Environmental and Economic Context}
Given the environmental and economic urgency of modern aviation, managing fuel consumption in the TMA is of utmost importance. The International Civil Aviation Organization (ICAO) has established ambitious global environmental goals, including continuous improvements in fuel efficiency and reducing the environmental footprint of aviation—directly supporting the United Nations' 2030 Agenda for Sustainable Development \cite{doc9750}. Inefficiencies in terminal areas significantly undermine these targets, as operations at lower altitudes burn fuel at much higher rates than during the cruise phase. Therefore, identifying, quantifying, and predicting these inefficiencies is crucial for mitigating the aviation industry's environmental impact and reducing operational costs.

\subsection{Performance-Based Management Framework}
To address these challenges and foster a more efficient global air transport network, ICAO advocates for performance-based management. This framework encourages stakeholders to engage in continuous performance benchmarking, guided at a strategic level by the Global Air Navigation Plan (GANP). In Brazil, the Department of Airspace Control (DECEA), supported by the ATM Performance Commission (CP-ATM), aligns with these global directives by defining operational metrics and conducting post-operational analyses of the Brazilian Airspace Control System (SISCEAB) \cite{mca10022}. By measuring specific outcomes, operators and regulators can share strategic goals and implement data-driven improvements.

\subsection{Machine Learning for Regression Tasks}
Machine learning has emerged as a powerful tool for solving complex regression problems across various domains. In the context of air traffic management, regression models can learn complex, non-linear relationships between operational variables and target metrics such as fuel consumption. Supervised learning approaches, particularly ensemble methods like gradient boosting, have demonstrated strong performance on structured tabular data by combining multiple weak learners to create a robust predictive model. These methods are particularly well-suited for capturing intricate interactions between traffic density, meteorological conditions, and operational constraints that characterize terminal area operations.
